% Local proof-prose fragment for the theta-star finite-atlas extension.
% This fragment is not a standalone paper and is not a main-paper update.

\section{Transport and theta-star invariance}

For a tree \(T\), let \(\operatorname{Rows}(T)\) be the eight equal-magnitude sign rows, and let \(\rho(T,s)\) be the zero-thickness row reach from the preceding section.  The proof spine proves \(t^*(T)\) by combining lower witnesses and upper caps over these eight rows; it does not need an exact first-contact formula for every non-attaining row.  The endpoint value \(\sqrt{3}\) means endpoint reached, not first contact.

\begin{lemma}[Eight-row bijection]
The signed transport rule
\[
  s'_{g(h)} = \det(R_g)\, \operatorname{orient}_g(h)\, s_h
\]
defines a bijection \(\operatorname{Rows}(T)\to\operatorname{Rows}(g(T))\).
\end{lemma}

\begin{proof}
The hinge map \(h \mapsto g(h)\) is a bijection on the three hinges.  For each hinge, the transported sign is obtained from the source sign by multiplication by one fixed element of \(\{+1,-1\}\).  Therefore the map is an invertible coordinate permutation followed by coordinatewise sign changes.  Such a map is a bijection on \(\{+1,-1\}^3\).
\end{proof}

\begin{lemma}[Row predicate preservation]
Let \(s\) be a row of \(T\), and let \(s'\) be its transported row.  For every admissible \(t\), the row \((T,s)\) is free, in contact, or obstructed exactly when \((g(T),s')\) is free, in contact, or obstructed after relabelling by the transported certificate data.  Consequently lower witnesses, upper caps, and exact row events are preserved under transport.
\end{lemma}

\begin{proof}
This is the finite-angle scaffold conjugacy lemma applied at the fixed value \(t\).  Pairwise relative transforms are preserved by piece relabelling and root re-gauge.  The finite SAT/contact predicates are rigid-motion invariant after the recorded axis-sign normalization.  Therefore the row status is preserved at that value of \(t\).  Since the same argument applies for every admissible \(t\), the row reach, lower-witness statements, and upper-cap statements are preserved.
\end{proof}

\begin{theorem}[Theta-star orbit invariance]
For every valid scaffold isometry \(g\),
\[
  t^*(T)=t^*(g(T)).
\]
The final theta-star class is also preserved: endpoint reached at \(\sqrt{3}\), positive jam at \(\sqrt{2}\), instant jam at \(0\), or the source-locked one-parameter wrapper class.
\end{theorem}

\begin{proof}
By the eight-row bijection, taking the maximum over the eight source rows is the same as taking the maximum over the eight transported target rows.  By row predicate preservation, transported lower witnesses and transported upper caps have the same values as their sources.  Hence the maximum certified row reach is unchanged.  The same argument preserves the semantic class attached to the event value, with \(\sqrt{3}\) interpreted as endpoint reached, not first contact.
\end{proof}

\begin{theorem}[S4 equal-magnitude theta-star finite-atlas theorem, local form]
In the historical zero-thickness S4 scaffold, over the finite atlas of 108 connected three-hinge trees and the eight equal-magnitude sign rows for each tree, the theta-star distribution certified by the proof spine is
\[
\begin{array}{rcl}
  40 &:& \sqrt{3} \text{ endpoint reached},\\
  8  &:& \sqrt{2} \text{ positive jam},\\
  60 &:& 0 \text{ instant jam}.
\end{array}
\]
Equivalently, the endpoint-reaching records split as \(36+4\), so the final class distribution is
\[
  36+4+8+60=108.
\]
The value \(\sqrt{2}\) corresponds to \(\theta = 2\arctan(\sqrt{2}) = \arccos(-1/3)\), the tetrahedral bond angle, i.e. the supplement of the regular tetrahedron dihedral angle.
\end{theorem}

\begin{proof}
The endpoint-free, positive-jam, and instant-jam classes have exact local theta-star certificates recorded by the representative table and source gates.  The wrapper class is recorded separately: its four records are source-locked one-parameter wrapper semantics inherited from public/local wrapper evidence, not a new broad motion certificate.  The finite-angle transport ledger maps the source records across the 108-tree atlas.  Root re-gauge replay proves that non-root-preserving records do not change pairwise geometry.  Endpoint-free witness transport records replace the historical candidate label for the full-open class.  The positive-jam source packages prove selected rows attaining \(\sqrt{2}\), and the positive-jam max-over-8 certificate proves the matching \(\sqrt{2}\) upper cap for every equal-magnitude row in those representatives.  The instant-jam representatives are closed by eight-row right-germ gates.  The theorem assembly and sabotage gates verify coverage, class counts, event values, and fail-closed behavior of the assembled proof spine.
\end{proof}
